\section{Data and replay}\label{app:data}

% ledger: M3, M4, M6, M7, Q22; source APPENDIX_DATA.md §1-§4, §6; verify/replay.py; viewer/README.md
\emph{Inputs.} The theorem is about one solid, \file{solid/r44_solid.json}: its $2{,}138$
vertices as exact rationals, its $4{,}272$ triangles, the $192$ features with their signed
coefficients, and the eight child poses. Two certificate files carry the finite data the proof
uses: \file{certificates/candidate_certificate.json} (the contact closure and profile, the
$44$-contact atlas, the $33$ first shells, the completions and conflicts, the parent census) and
\file{certificates/companion_collision_certificate.json} (the $5{,}273$ rejection records with
their $299{,}975$ collision witnesses, the boxes themselves in a companion file). Their sha256
digests are in \cref{tab:hashes}. The replay entrypoint checks that the four canonical inputs are
byte-identical to the packet copies and checks the fixed solid digest before invoking the
checkers.

\begin{table}[ht]
  \centering
  \caption{Input hashes and tool versions, computed at the state of the repository described in
  this paper (FIGURES.md, Table 6).}
  \label{tab:hashes}
  \small
  % generated by paper/scripts/appendix_data_table.py; sha256 computed live; Lean versions at d90313a717
\begin{tabular}{L{0.42\linewidth} L{0.5\linewidth}}
\toprule
artifact & sha256 / version \\
\midrule
\texttt{solid/\allowbreak{}\allowbreak{}r44\_\allowbreak{}solid.json} & \texttt{f320d7a0c2d784a3\dots{}42f0ed55} \\
\texttt{solid/\allowbreak{}\allowbreak{}native\_\allowbreak{}panels.csv} & \texttt{b9eec9fd834abba4\dots{}d3a9ed2d} \\
\texttt{certificates/\allowbreak{}\allowbreak{}candidate\_\allowbreak{}certificate.json} & \texttt{44e9b3f6048497de\dots{}c5de5ff2} \\
\texttt{certificates/\allowbreak{}\allowbreak{}companion\_\allowbreak{}collision\_\allowbreak{}certificate.json} & \texttt{36e89e2fe81f4788\dots{}1b068b08} \\
\texttt{certificates/\allowbreak{}\allowbreak{}collision\_\allowbreak{}core\_\allowbreak{}witnesses.jsonl.gz} & \texttt{ae64a5283400c61e\dots{}161f8a45} \\
Lean toolchain & \texttt{leanprover/\allowbreak{}lean4:v4.31.0} \\
Mathlib commit & \texttt{fabf563a7c95} \\
batteries / aesop / Qq / proofwidgets & \texttt{fa08db58b30e, e3cb2f741431, f46324995fca, 24b0d9dc081c} \\
\texttt{verify/\allowbreak{}\allowbreak{}replay.py} (Python standard library) & PASS, 27.0 s \\
Python (replay and tables) & 3.12.3 \\
\bottomrule
\end{tabular}

\end{table}

% ledger: M3, M4, M5, M6; source APPENDIX_DATA.md §2-§4; paper/review/REPLAY_FROM_PAPER.md
\emph{Replay in one minute.} From the repository root,
\begin{verbatim}
python3 verify/replay.py
\end{verbatim}
runs, with the Python standard library only and no network, the checker of the registered
theorem (the first implementation), the alignment checker and the integer-only replay of the
collision boxes (the second implementation), the mesh check \file{verify/boundary_sphere.py},
the tube-formula controls, \file{verify/substitution_modular_coincidence.py} and the
block-language closure \file{verify/seed_language_closure.py}; it asserts the
input digests first and writes \file{verify/replay_report.json}. In the recorded clean-clone run
at the pinned commit (\file{paper/review/REPLAY_FROM_PAPER.md}, 2026-09-10) the replay
completed in about half a minute ($28.316$
seconds wall time) on the machine that ran the Lean builds; its hardware is not otherwise
specified; hardware specifications and Lean build time and peak memory are not reported for this
version. Three
further commands complete the README's one-minute list:
\begin{verbatim}
(cd lean/R44 && python3 gen/extract.py --check)
python3 -m unittest discover simulations/tests
cd verify/packets/r44_unrestricted_alignment
python3 src/test_mutations.py
\end{verbatim}
The first checks that the Lean literals generated from the JSON inputs match those inputs. The
unit-test command reported $10$ tests with six skipped by optional-dependency guards; it does not
establish that those six tests passed. The last command rejects six corrupted inputs (a missing triangle, a reversed triangle, a moved apex, a
missing rejection witness, a wrong companion quantifier, a missing registered contact) and states
its own scope: it checks that malformed inputs are rejected and does not validate the continuous
proof.

\emph{Paper-lane scripts.} The standard-library scripts cited in the text as \tier{T2} live in
\file{paper/scripts/}: the exact closure of the block language and the fixed seeds
(\file{seed_language_closure.py}, \cref{sec:limitperiodic}), the planar areas of \cref{app:asymmetry}
(\file{planar_areas.py}) and the non-convexity witness; each reads the repository data and
prints \texttt{STATUS: PASS}.

\emph{Figures.} \Cref{fig:tile,fig:panel} are drawn by \file{figures/print-preview/build\_preview.py}
from \file{solid/r44\_solid.json} with the exaggerations stated in their captions: the carrier and
the $192$ feature centres keep their coordinates, every feature base is widened $\times 5$ about
its centre and every signed height is multiplied by $80$, so polarities and relative heights
are preserved and nothing is added, omitted or moved. The script checks all $192$ feature records
against the mesh and the native panel table in rational arithmetic, and a second script reads the
display mesh back and reconstructs its $768$ feature triangles independently; the outputs and their
digests are in \file{figures/print-preview/}. \Cref{fig:cluster,fig:patch} are the renders of
\file{figures/render\_r44.py}. No figure is verification.

\emph{Exact arithmetic.} Every certificate is integer or rational data; the only decimal literals
in the certificates are two timing fields. The checkers use Python integers and
\texttt{fractions.Fraction}; the only floating-point conversion in the repository's verification
code is the export of the mesh to a rendering format. The verification scripts named in the ``Paper-lane scripts'' paragraph
also use exact rationals only; the planar-area checker bounds each non-axis facet's
area $\sqrt q/2$ by $(q+m^2)/(4m)$ with an integer-square-root witness $m$, exactly. Collision
boxes are stored in integer units of $1/400$. The Lean theorems decide integer and rational
literals.

\emph{Lean.} The development is \file{lean/R44}, built with \texttt{lake build} under the
toolchain of \cref{tab:hashes} (the build is capped at four parallel jobs and the Mathlib cache is
pre-seeded); \file{lean/R44/scripts/controls.sh} runs the positive build, then every must-fail
file against its expected diagnostic, the scope regressions and the admission count, in one
command. The axiom lines of \cref{tab:axioms} are produced by the build itself
(\file{build_axioms.log}).

% ledger: Q22, M7, O12; source viewer/README.md; figures/print-preview/README.md
\emph{Viewer.} \file{viewer/index.html} shows the solid, the eight-child cluster and the
$64$-, $512$- and $4{,}096$-copy patches in the browser without a server. It
is generated from the canonical rational data with the source hashes recorded; its features are
drawn exaggerated by default, with a switch to true heights; colours are annotation. It also
offers a period test on finite patches, which, as its own caption says, proves nothing.


\emph{Code.} Repository: \url{https://github.com/ioannist/six-birds-tiles}. The Lean statements,
hypothesis record and axiom lists reproduced in this paper refer to commit
\texttt{d90313a717}, and the recorded clean-clone replay was run at that commit (an earlier
recorded run at \texttt{bdc6d3b957}, before the block-language closure joined the replay, gave
the same digests and outcomes). The intended paper-release tag is \texttt{paper-v1};
tag creation and licensing remain pending owner actions.
